EV化、都市物流、last-mile delivery、充電インフラ制約、depot/grid制約は、配送計画問題を単なるVehicle Routing Problem (VRP)から、容量制約、時間窓、充電状態、再最適化を含むCVRP/EVRPへ拡張している。
一方、量子最適化研究ではQAOA、VQA、QUBO/Ising定式化、HOBO、minimal encodingなどを用いたVRP/CVRP/EVRPへの適用が進められている。

しかし、既存研究で扱われる問題規模や回路資源は、多くの場合、理論定義、小規模simulation、小規模NISQ実験、あるいはresource estimationの段階にとどまる。
したがって、本研究では量子計算が古典手法に対して実用的に優位であると主張するのではなく、既存研究がどの技術段階にあり、将来のEV・物流システムが要求する社会段階との間にどのようなギャップが残存しているかを評価する。

本研究の中心問いは次である。

\begin{quote}
交通分野における量子技術開発と社会的要請の関係は、どのように評価可能か。
\end{quote}

この問いに答えるため、本研究では既存量子VRP/CVRP/EVRP研究からwidth、depth、gate count、shots、feasibility、workflowに関する情報を抽出し、社会段階と技術段階を接続するgap matrixとして整理する。
